\documentclass[12pt,a4paper]{article}
\usepackage{ucebnice}
\usetikzlibrary{fit,shapes}

\title{Funkce}
\author{Ondřej Škubala}
\date{\today}

\begin{document}

\maketitle
\newpage


\section{Funkce}


\begin{quotebox}
{„Funkce popisují svět. Všechno kolem nás lze vyjádřit pomocí funkcí – zvuk, který slyšíme, světlo, které dopadá na naše oči, i vzorce, které tvoří naše data.  
Každá oblast matematiky se zabývá jinými druhy funkcí: středoškolská matematika zkoumá polynomy jedné proměnné, analýza se věnuje spojitým funkcím… a tak to pokračuje dál.  

FUNKCE POPISUJÍ SVĚT!“}
{Thomas A. Garrity, \textit{\href{https://www.youtube.com/watch?v=BWZTlfrneD8}{Functions describe the world}}}
\end{quotebox}

\begin{historicalbox}
Pojem \textbf{funkce}, který dnes používáme s naprostou samozřejmostí, má překvapivě dlouhou a spletitou historii.  
Lidé už od starověku pozorovali, že určité veličiny se mění v závislosti na jiných – například že délka stínu souvisí s výškou Slunce, že rychlost pohybu ovlivňuje dráhu, nebo že tlak závisí na objemu.  
Matematika těchto vztahů se však tehdy nevyjadřovala symbolicky, ale \emph{geometricky}: pomocí křivek, úhlů a podobností.  

\medskip
\textbf{Antika a středověk: geometrie jako jazyk závislostí}

Už \textbf{Eukleidés} (3. století př. n. l.) a jeho následovníci vyjadřovali vztahy mezi veličinami geometricky.  
Křivka byla obrazem závislosti – její tvar prozrazoval, jak se jedna veličina mění vůči druhé.  
Později \textbf{Apollónios z Pergy} popsal kuželosečky (elipsu, parabolu a hyperbolu), které se staly prvním „modelem funkcí“:  
popisovaly fyzikální i astronomické zákonitosti, aniž by někdo použil algebraický zápis.  

Ve středověku se s těmito geometrickými představami dále pracovalo, ale až s nástupem \textbf{analytické geometrie} v 17. století došlo k zásadnímu zlomu – matematika dostala možnost převést geometrické vztahy do číselné a symbolické podoby.

\medskip
\textbf{Galileo, Torricelli a vznik funkčního myšlení}

Další vývoj funkčního myšlení je spojen s pracemi \textbf{Galilea Galileiho} (1564–1642) a \textbf{Evangelisty Torricelliho} (1608–1647).  
Ve svých výzkumech pohybu a fyzikálních jevů sledovali vztahy mezi veličinami pomocí křivek – například dráhu tělesa při volném pádu nebo tvar proudící kapaliny.  
Jejich přístup byl stále geometrický, ale už v sobě nesl myšlenku závislosti jedné veličiny na druhé – tedy to, co dnes označujeme jako funkční vztah.

\medskip
\textbf{René Descartes a zrod analytické geometrie}

Francouzský filozof a matematik \textbf{René Descartes} (1596–1650) poprvé systematicky zavedl souřadnicový systém, v němž každému bodu v rovině odpovídá číselná dvojice \((x, y)\).  
Tím umožnil popisovat křivky pomocí rovnic a otevřel cestu k symbolickému vyjádření závislostí – tedy k pojmu, který dnes nazýváme funkcí.  
Křivka tak přestala být pouze geometrickým objektem a stala se obrazem číselné závislosti mezi veličinami.  
Descartes tím doslova sjednotil aritmetiku s geometrií.

\newpage
\textbf{Leibniz, Bernoulli a první definice funkce}

Samotné slovo \emph{funkce} (\textit{functio}) se poprvé objevilo roku 1673 v dílech \textbf{Gottfrieda Wilhelma Leibnize} (1646–1716).  
Původně jím označoval geometrické vlastnosti křivek, například tečny či body zakřivení.  
V roce 1679 dal první definici pojmu funkce \textbf{Johann Bernoulli} (1667–1748), který pojem dále rozvíjel spolu s Leibnizem a \textbf{Isaacem Newtonem} (1643–1727) – zakladateli diferenciálního a integrálního počtu.  
Funkce se tehdy chápala jako veličina, která „vykonává činnost“ – tedy vyjadřuje vztah mezi proměnnými.

\medskip
\textbf{Euler, Fourier a Dirichlet – zrod moderního pojetí}

Zásadní krok k dnešnímu chápání udělal švýcarský matematik \textbf{Leonhard Euler} (1707–1783).  
Ve svém díle \textit{Introductio in analysin infinitorum} (1748) definoval funkci jako:
\begin{quote}
„Funkcí proměnné veličiny rozumíme jakýkoli analytický výraz složený libovolným způsobem z této veličiny a z konstant.“
\end{quote}
V Eulerově době se tedy za funkci považoval každý \emph{vzorec}, který dává smysl – například polynom, zlomek či odmocnina.  
Euler také jako první běžně používal zápis \(f(x)\) a systematicky rozlišoval \textit{nezávislou proměnnou} a \textit{závislou veličinu}.  

Na přelomu 18. a 19. století pak \textbf{Joseph Fourier} (1768–1830) a \textbf{Peter Gustav Lejeune Dirichlet} (1805–1859) rozšířili pojem funkce i na případy, kdy ji nelze vyjádřit jednoduchým vzorcem.  
Dirichlet formuloval novou definici:
\begin{quote}
„Funkce \(y\) proměnné \(x\) je zákon, podle něhož každé hodnotě \(x\) z určitého oboru přísluší právě jedna hodnota \(y\).“
\end{quote}
Tím byla funkce osvobozena od nutnosti být vyjádřena analyticky – stačilo, že mezi hodnotami existuje jednoznačné přiřazení.  

\medskip
\textbf{Cauchy, Weierstrass a formální analýza}

Na Dirichletovy myšlenky navázali \textbf{Cauchy}, \textbf{Weierstrass} a další, kteří zavedli pojem \emph{limity}, \emph{spojitosti} a \emph{derivace} ve formální podobě.  
Matematika se tím stala přísnější a pojem funkce dostal pevný logický základ.

\medskip
\textbf{20. století – funkce jako zobrazení mezi množinami}

Ve 20. století se pojem funkce dále zobecnil.  
Díky rozvoji \textbf{teorie množin} (Cantor, Bolzano) a \textbf{matematické logiky} se funkce začala chápat jako \emph{zobrazení} mezi množinami, tedy jako soubor uspořádaných dvojic \((x, y)\) s vlastností jednoznačnosti.  
Tento přístup sjednotil funkce reálné, komplexní, vektorové i abstraktní.  
Funkce se staly základním pojmem nejen v analýze, ale i v \textbf{algebře}, \textbf{geometrii}, \textbf{pravděpodobnosti} a \textbf{informatice}.  

\medskip
\textbf{Současnost – funkce jako univerzální princip}

V moderním pojetí lze funkci chápat nejen jako matematický objekt, ale jako \emph{obecný princip popisu závislostí}.  
Funkce jsou základem programování (každá funkce v kódu představuje přiřazení vstup–výstup), modelování dat, fyzikálních simulací i neuronových sítí.  
Myšlenka, že „každý vstup má svůj výstup“, je natolik univerzální, že prostupuje celou vědou, technikou i logikou.
\end{historicalbox}



\begin{notebox}
\textbf{Zajímavost:}  
Zatímco dnes zapisujeme funkce jako \(y = f(x)\), Euler často používal označení \(\varphi(x)\) nebo jednoduše „funkce \(x\)“.  
Teprve postupem času se zápis \(f(x)\) stal standardem, který přetrval dodnes.
\end{notebox}


Jak je z citátu pana \textbf{Thomase A. Garrityho} zřejmé, pojem \important{funkce} představuje nesmírně široké téma, se kterým se nyní podrobněji seznámíme.  
Funkce tvoří jádro celé matematiky – nejsou jen jednou z kapitol, ale základním jazykem, pomocí něhož dokážeme popisovat děje, závislosti i změny.  

S pojmem funkce se nesetkáme pouze nyní, ale bude nás provázet po zbytek studia:  
v dalších ročnících při práci s \textbf{lineárními, kvadratickými, exponenciálními a logaritmickými funkcemi},  
později také u \textbf{posloupností, limit, derivací} a mnoha dalších matematických oblastí.  
Proto je pojem \important{funkce} naprosto \textbf{klíčový} a jeho pochopení představuje důležitý krok k hlubšímu porozumění matematice jako celku.

Funkci si můžeme představit jako \textbf{pravidlo}, které převádí jeden typ vstupu na odpovídající výstup.  
Toto pravidlo může být vyjádřeno matematickým vzorcem, logickou podmínkou, algoritmem nebo dokonce fyzikálním či biologickým procesem.

\begin{examplebox}
\textbf{Příklady různých typů funkcí:}
\begin{itemize}
  \item \textbf{Matematická funkce:}  
  Například \(f(x) = 2x + 3\) přiřadí každému číslu \(x\) právě jedno číslo \(y\).  
  Je to klasický příklad lineární závislosti.

  \item \textbf{Speciální funkce: Dirichletova funkce}  
  \[
  f(x) = 
  \begin{cases}
  1, & \text{je-li } x \in \mathbb{Q},\\
  0, & \text{je-li } x \notin \mathbb{Q}.
  \end{cases}
  \]
  Tato funkce rozlišuje, zda je číslo racionální nebo iracionální – i to je platná funkce, přestože její graf je „roztříštěný“ po celé číselné ose.

  \item \textbf{Informatická funkce:}  
  V programování má podobný princip například příkaz:
  \[
  \texttt{if (x > 0) return 1; else return 0;}
  \]
  který každému vstupu \(x\) přiřadí 1 nebo 0 – opět jednoznačně.

  \item \textbf{Fyzikální funkce:}  
  Tlak vzduchu \(p(h)\) závisí na nadmořské výšce \(h\); se stoupající výškou tlak klesá.  
  I to je funkční závislost mezi dvěma veličinami.

  \item \textbf{Biologická funkce:}  
  Plíce lze chápat jako „funkci“, která každé neokysličené krvi přiřadí krev okysličenou.  
  I v biologii tedy najdeme procesy, které přesně odpovídají pojmu funkce – mají vstup, mechanismus a jednoznačný výstup.

  \item \textbf{Ekonomická funkce:}  
  Například funkce úroku v bance:
  \[
  f(p) = p \cdot (1 + r)^t,
  \]
  kde \(p\) je počáteční vklad, \(r\) úroková míra a \(t\) počet let.  
  Každému vkladu jednoznačně přiřadí jeho budoucí hodnotu.  
  (Zde se poprvé setkáváme s exponenciálním růstem.)

  \item \textbf{Sociologická funkce:}  
  Počet obyvatel města \(P(t)\) jako funkce času \(t\) – může růst, klesat, nebo se stabilizovat.  
  Takový model umožňuje sledovat demografický vývoj či plánovat kapacitu infrastruktury.

  \item \textbf{Psychologická funkce:}  
  Reakční čas člověka \(R(s)\) může být funkcí síly podnětu \(s\):  
  při malém podnětu reagujeme pomalu, při silném rychleji – až do určitého limitu.  
  I zde vidíme funkční závislost, byť biologicky a experimentálně měřenou.
\end{itemize}
\end{examplebox}

\begin{notebox}
Pojem funkce přesahuje hranice matematiky – vyskytuje se všude tam, kde \textbf{jedné veličině odpovídá přesně jedna jiná}.  
Ať už jde o úrok v bance, tlak vzduchu nebo reakci organismu, princip je vždy stejný:  
\textit{vstup → pravidlo → výstup.}
\end{notebox}



\section{Co je funkce}

\important{Funkce} patří k \textbf{nejzákladnějším pojmům} celé matematiky.  
Umožňují popsat závislost jedné veličiny na druhé – například jak se mění teplota v průběhu dne, jak roste cena zboží při různé hmotnosti, nebo jak se mění výška tělesa při volném pádu v čase.  
Matematicky lze tuto závislost popsat několika různými, ale vzájemně ekvivalentními způsoby.  
Níže jsou uvedeny tři nejčastější formy, s nimiž se setkáváme.

\begin{definitionbox}
    \textbf{1. Funkce jako zobrazení mezi množinami}

    Funkce \( f \) je \emph{zobrazení} z množiny \( X \) do množiny \( Y \), 
    při němž každému prvku \( x \in X \) je přiřazen právě jeden prvek \( y \in Y \).  
    Zapisujeme:
    \[
    f: X \rightarrow Y, \quad x \mapsto f(x).
    \]
    Množina \( X \) se nazývá \textbf{definiční obor} (domain), množina \( Y \) \textbf{obor hodnot} (codomain).
\end{definitionbox}

Tento pohled na funkci bývá nejpoužívanější ve školní matematice, 
protože je přirozený a snadno se znázorňuje.  
Můžeme si ho představit jako „přiřazovací stroj“ – do množiny \( X \) 
vkládáme vstupy a funkce nám pro každý z nich vrátí výstup v množině \( Y \).  
Příklad: pokud \( X \) představuje množinu lidí a \( Y \) množinu jejich dat narození, 
pak funkce „datum narození“ každému člověku přiřadí právě jedno datum.

\smallskip

\begin{definitionbox}
    \textbf{2. Funkce jako množina uspořádaných dvojic}

    Funkce může být chápána také jako množina všech dvojic \((x, y)\), kde \(x \in X\) a \(y \in Y\), tedy jako podmnožina kartézského součinu:
    \[
    f \subseteq X \times Y.
    \]
    Musí přitom platit, že pro každé \(x \in X\) existuje právě jedno \(y \in Y\), pro které \((x,y) \in f.\)
\end{definitionbox}

Tento přístup už zdůrazňuje množinový základ celé matematiky.  
Místo toho, abychom si funkci představovali jako proces, 
vnímáme ji spíš jako \textit{sadu konkrétních přiřazení}.  
Každá dvojice \((x, y)\) vyjadřuje, že hodnota \(y\) „patří“ k argumentu \(x\).  
Například pro funkci \(f(x) = 2x + 1\) je tato množina všech dvojic rovna \(\{(x, 2x+1) \mid x \in \mathbb{R}\}\).  
Takové pojetí je základem pro další práci s funkcemi jako s množinovými objekty (sjednocování, průniky, skládání atd.).

\smallskip

\begin{definitionbox}
    \textbf{3. Funkce jako speciální binární relace}

    Funkci lze chápat i jako \emph{binární relaci} mezi množinami \(X\) a \(Y\), která splňuje podmínku jednoznačnosti:
    \[
    \forall x \in X, \; \exists! \, y \in Y : (x, y) \in f.
    \]
    Zapisujeme-li funkci tímto způsobem, zdůrazňujeme její logický charakter – jako vztah mezi prvky dvou množin.
\end{definitionbox}

Tento způsob zápisu bývá nejužitečnější, pokud se na funkce díváme z logického 
nebo teoretického hlediska.  
Funkce se zde chápe jako „zvláštní případ relace“, 
která má jednu specifickou vlastnost – ke každému prvku \(x\) 
existuje právě jeden prvek \(y\).  
Tento přístup používají pokročilejší oblasti matematiky, 
jako je teorie množin, logika nebo informatika, 
kde je potřeba mít přesně definováno, jak se prvky různých množin propojují.

\medskip
Každý z uvedených přístupů tedy vychází ze stejné myšlenky – že funkce \emph{jednoznačně přiřazuje výstup každému vstupu} – jen ji popisuje jiným jazykem.  
První pohled je názorný, druhý množinový a třetí logicko-formální.  
Všechny tři však vyjadřují tentýž pojem, který se stává základem celé matematiky.

\begin{notebox}
    \textbf{Představa funkce jako „black-boxu“}

    Nejjednoší způsob jak si představit funkci je si ji představit jako černou skříňku, do které vhodíme číslo \(x\).  
    Uvnitř sedí pomyslný skřítek, který podle určitého pravidla něco provede
     – možná sečte, možná umocní nebo odečte – 
     a nakonec nám ze skříňky vyhodí číslo \(y\).  
    Není důležité vědět, co přesně skřítek dělá, ale platí, 
    že na tentýž vstup vždy dostaneme tentýž výstup.  

    \medskip
    Takto si můžeme snadno představit zápis \( f(x) = 2x + 1 \):  
    do skříňky vhodíme číslo \(x\), skřítek ho zdvojnásobí, přičte jedničku – a ven pošle výsledek \(y\).
\end{notebox}


\section{Definiční obor a obor hodnot}

Abychom mohli přesně určit, kdy má funkce smysl a jakých hodnot může nabývat, 
zavádíme dva důležité pojmy – \important{definiční obor} a \important{obor hodnot}.  
Definiční obor říká, \emph{které vstupy} jsou pro funkci přípustné, 
zatímco obor hodnot určuje, \emph{jaké výstupy} může funkce přijímat.  
Ne vždy totiž platí, že funkce využívá všechny prvky množiny \(Y\), 
do které zobrazuje – někdy určité hodnoty vůbec nenastanou.

\subsubsection*{Definiční obor}

Definiční obor funkce je množina všech hodnot proměnné \(x\), pro které má zápis \(f(x)\) smysl.  
Jinými slovy, jsou to všechna čísla, která \emph{můžeme dosadit za \(x\)}, aniž by došlo k nesmyslné situaci – například dělení nulou, odmocnění záporného čísla (v reálných číslech), logaritmus záporného čísla apod.  

Z hlediska grafu představuje definiční obor všechny \textbf{projekce bodů grafu funkce na osu \(x\)} – tedy všechny hodnoty \(x\), pro které na grafu existuje odpovídající bod.

S pojmem definiční obor jsme se už setkali při \important{algebraických výrazech}, kde jsme například u \emph{lomených výrazů} určovali, pro která čísla má daný výraz smysl.  
Funkce tyto myšlenky pouze zobecňuje:  
zjišťujeme, pro která \(x\) má výraz \(f(x)\) význam a pro která nikoli.  

\begin{examplebox}
    Pokud má například funkce tvar
    \[
    f(x) = \frac{2x + 1}{x - 3},
    \]
    pak musíme vyloučit hodnotu \(x = 3\), protože jmenovatel by se rovnal nule.  
    Definiční obor je tedy \(D(f) = \mathbb{R} \setminus \{3\}\).
\end{examplebox}



\subsubsection*{Obor hodnot}

Obor hodnot je množina všech výstupních hodnot \(y\), kterých funkce může skutečně nabýt při všech přípustných \(x\).  
Jinými slovy, obor hodnot odpovídá tomu, \emph{kam až funkce „dosáhne“} – které \(y\) se na grafu skutečně vyskytují.  
Na grafu tvoří obor hodnot všechny \textbf{projekce bodů grafu na osu \(y\)}.

Obor hodnot tedy neříká, \emph{co by funkce mohla mít}, ale \emph{co skutečně má}.  
Zatímco definiční obor určuje, co smíme dosadit za \(x\), obor hodnot určuje, které výsledky z takového dosazování mohou vzniknout.  
Při práci s funkcemi je obor hodnot často možné zjistit ze samotného předpisu, z jeho tvaru nebo z vlastností známých funkcí (např. že druhá mocnina je vždy nezáporná).

\begin{examplebox}
Uvažujme funkci
\[
f(x) = \frac{2x + 1}{x - 3}.
\]
Z předchozího víme, že její definiční obor je \(D(f) = \mathbb{R} \setminus \{3\}\).

Abychom určili obor hodnot, ptáme se: jakých všech hodnot může výraz \(\frac{2x + 1}{x - 3}\) nabývat?  
Zjistíme, že žádné dosazení nedá výsledek \(y = 2\).  
To lze ověřit algebraicky:
\begin{center}
    \[
    \begin{array}{rl}
        y &= \dfrac{2x + 1}{x - 3} \\[6pt]
        yx - 3y &= 2x + 1 \\[6pt]
        yx - 2x &= 3y + 1 \\[6pt]
        x(y - 2) &= 3y + 1 \\[6pt]
        x &= \dfrac{3y + 1}{y - 2}
    \end{array}
    \]
\end{center}



Rovnice má řešení pro všechna \(y\) kromě \(y = 2\), protože při \(y=2\) by se vyskytlo dělení nulou.  
Tedy obor hodnot je 
\[
H(f) = \mathbb{R} \setminus \{2\}.
\]
\end{examplebox}

Z grafického hlediska má tato funkce tzv. \emph{asymptoty} – přímky, ke kterým se graf blíží, ale nikdy jich nedosáhne.  
Svislá asymptota \(x=3\) odpovídá vyloučené hodnotě z definičního oboru,  
vodorovná asymptota \(y=2\) pak vyloučené hodnotě z oboru hodnot.

\begin{notebox}
\textbf{Shrnutí:}  
    \begin{itemize}
      \item \textbf{Definiční obor} – co můžeme do funkce „vložit“ (vstupy).  
      \item \textbf{Obor hodnot} – co z funkce může „vypadnout“ (výstupy).  
      \item Na grafu: definiční obor odpovídá projekci bodů na osu \(x\),  
            obor hodnot projekci bodů na osu \(y\).
    \end{itemize}
\end{notebox}


\subsection{Způsoby zadání funkce}

Funkci můžeme popsat různými způsoby – podle toho, co chceme zdůraznit.  
Někdy nám stačí slovní vyjádření, jindy konkrétní tabulka nebo přesný vzorec.  
Všechny tyto způsoby jsou rovnocenné, pouze ukazují tentýž vztah mezi proměnnými z různých pohledů.

\medskip
\textbf{1. Slovní popis funkce}

Nejjednodušší a nejpřirozenější způsob, jak funkci popsat, je pomocí slov.  
Funkce je v takovém případě vyjádřena jako vztah mezi veličinami, který známe z praxe:
\begin{quote}
    \emph{
        „Každému dni přiřadíme průměrnou denní teplotu.“  \\
        „Každému člověku přiřadíme jeho datum narození.“  \\
        „Každé délce strany přiřadíme obsah čtverce se stejnou stranou.“ \\
        } 
\end{quote}

Takový zápis bývá vhodný tam, kde jde spíše o smysl závislosti než o její přesný výpočet.

\medskip
\textbf{2. Tabulkový zápis}

Funkci lze vyjádřit také \emph{tabulkou hodnot}, kde v jednom sloupci (nebo řádku) uvádíme vstupy – hodnoty \(x\), a ve druhém odpovídající výstupy \(f(x)\).  
Tabulka názorně ukazuje, jak se výstup mění podle vstupu.  
Z takové tabulky můžeme snadno přejít ke grafickému znázornění.

\begin{center}
    \begin{minipage}{0.45\textwidth}
        \centering
        \arrayrulecolor{tableborder}
        \begin{tabular}{|c|c|}
            \hline
            \rowcolor{tableheaderdark} \(x\) & \(f(x) = 2x + 1\) \\
            \hline
            $-2$ & $-3$ \\
            \hline
            $-1$ & $-1$ \\
            \hline
            $0$ & $1$ \\
            \hline
            $1$ & $3$ \\
            \hline
            $2$ & $5$ \\
            \hline
        \end{tabular}
        \arrayrulecolor{black}
        \captionof*{table}{\small Vertikální tabulka hodnot}
    \end{minipage}
            \hspace{1cm}
    \begin{minipage}{0.45\textwidth}
        \centering
        \arrayrulecolor{tableborder}
        \begin{tabular}{|c|c|c|c|c|c|}
            \hline
            \rowcolor{tableheaderdark} \(x\) & $-2$ & $-1$ & $0$ & $1$ & $2$ \\
            \hline
            \(f(x)=2x+1\) & $-3$ & $-1$ & $1$ & $3$ & $5$ \\
            \hline
        \end{tabular}
    \arrayrulecolor{black}
    \captionof*{table}{\small Horizontální tabulka hodnot}
    \end{minipage}
\end{center}


Obě tabulky vyjadřují stejný vztah mezi vstupy a výstupy – liší se pouze orientací.  
V některých případech je vertikální zápis přehlednější (např. pro větší množství hodnot), jindy naopak horizontální (např. pro snadné porovnání trendu).


\medskip
\textbf{3. Předpis (analytické vyjádření)}

Nejběžnější a nejpřesnější způsob zápisu funkce je \emph{předpisem} neboli vzorcem.  
Ten udává pravidlo, podle něhož se každému \(x\) přiřadí hodnota \(f(x)\). 
S takovým vzorcem jsme se již setkali, když jsme určovali definiční obor a obor hodnot $y=\frac{2x+1}{x-3}$
Dalšími takovými předpisy jsou apříklad:

\[
f(x) = 2x + 1, \quad g(x) = x^2 - 4x + 3, \quad h(x) = \frac{1}{x - 2}.
\]

Takový zápis umožňuje provádět výpočty, určovat vlastnosti funkce, její graf i právě již zmíněný definiční obor a obor hodnot.

\medskip
\textbf{4. Grafické znázornění}

Graf funkce je množina všech bodů \((x, y)\), pro které platí \(y = f(x)\).  
Každému bodu na grafu odpovídá dvojice hodnot – vstupní a výstupní.  
Grafické zobrazení umožňuje rychle odhadnout chování funkce: kde roste, kde klesá, zda je kladná nebo záporná, a kde má průsečíky s osami.

\begin{center}
  \begin{axisgraph}[
      grid=none,
      xmin=-2, xmax=3,
      ymin=-3, ymax=3,
      title={Graf přímky $y = 2x - 1$}
  ]
      \addplot[domain=-4:4, thick, blue]{2*x - 1};
  \end{axisgraph}
\end{center}



Graf je tedy vizuálním doplněním předpisu – ukazuje tutéž informaci, ale názorněji.  
V praxi se všechny tyto způsoby často kombinují: slovní popis doplníme vzorcem, k němu vytvoříme tabulku a z ní pak graf.



%Na obrázku níže je znázorněno zobrazení mezi množinami \(X\) a \(Y\).  
%Každému prvku z množiny \(X\) odpovídá právě jeden prvek v množině \(Y\).
%
%\begin{center}
%    \begin{tikzpicture} 
%        [
%            >=stealth, bullet/.style={fill=black, circle,minimum width=1pt,inner sep=1pt}, 
%            projection/.style={->,thick,shorten <=2pt,shorten >=2pt}, 
%            every fit/.style={ellipse,draw,inner sep=0pt} 
%        ] 
%
%        \foreach \y/\l in {1/d,2/c/,3/b,4/a} 
%        \node[bullet,label=left:$\l$] (a\y) at (0,\y) {}; 
%        
%        \foreach \y/\l in {1/4,2/3,3/2,4/1} 
%        \node[bullet,label=right:$\l$] (b\y) at (4,\y) {}; 
%        
%        \node[draw,fit=(a1) (a2) (a3) (a4),minimum width=2cm] {} ; 
%        \node[draw,fit=(b1) (b2) (b3) (b4),minimum width=2cm] {} ; 
%        
%        \draw[projection] (a1) -- (b4); \draw[projection] (a2) -- (b2); 
%        \draw[projection] (a3) -- (b1); \draw[projection] (a4) -- (b3); 
%    
%    \end{tikzpicture} 
%\end{center}


\newpage
\section{Vlastnosti funkcí}

Kromě samotného předpisu a grafu nás u funkcí zajímá i to, \textbf{jak se chovají}.  
Zda rostou nebo klesají, jestli mají největší nebo nejmenší hodnotu, zda jsou „symetrické“, nebo zda se jejich chování opakuje v určitých intervalech.  
Tyto charakteristiky označujeme jako \important{vlastnosti funkcí}.  

Níže jsou uvedeny nejdůležitější z nich, se kterými se v průběhu studia matematiky budeme často setkávat.



% --- D E F I N I Č N Í   O B O R   A   O B O R   H O D N O T ---
\subsection*{Definiční obor a obor hodnot}

Než přejdeme k dalším vlastnostem funkcí, připomeňme si dvě základní pojmy,
které jsme již probrali výš, avšak se bez nich neobejdeme.

\begin{definitionbox}
\textbf{Definiční obor funkce} (označujeme \(D(f)\)) je množina všech čísel \(x\), pro která má zápis \(f(x)\) smysl.  
Zjednodušeně řečeno: \textit{všechna čísla, která můžeme do funkce „vložit“ jako vstup.}
\end{definitionbox}

\begin{definitionbox}
\textbf{Obor hodnot funkce} (označujeme \(H(f)\)) je množina všech čísel \(y\), kterých funkce může nabývat.  
Jinými slovy: \textit{všechna čísla, která funkce může „vyhodit“ jako výsledek.}
\end{definitionbox}

\begin{examplebox}
Například pro funkci:
\[
f(x) = \frac{1}{x - 2}
\]

nesmí být jmenovatel nulový, takže \(x \neq 2\).  
Definiční obor je \(D(f) = \mathbb{R} \setminus \{2\}\).
A zároveň nikdy nedostaneme \(y = 0\) (protože tento zlomek nemůže být roven nule).  
Tedy \(H(f) = \mathbb{R} \setminus \{0\}\). A tyto osy nazýváme \important{asymptoty}.
\end{examplebox}

\begin{center}
  \begin{axisgraph}[grid=none, xmin=-4, xmax=6, ymin=-4, ymax=4, title={Graf funkce: \(f(x) = \dfrac{1}{x-2}\)}]
      \addplot[domain=-4:1.8, thick, blue, samples=100]{1/(x-2)};
      \addplot[domain=2.2:6, thick, blue, samples=100]{1/(x-2)};
      \addplot[dashed, gray] coordinates {(2,-4) (2,4)};
      \addplot[dashed, gray] coordinates {(-4,0) (6,0)};
  \end{axisgraph}
\end{center}


% --- S U D Á   A   L I C H Á   F U N K C E ---
\subsection*{Sudá a lichá funkce}

Některé funkce mají zajímavou \textbf{symetrii} – jejich grafy se „opakují“ určitým způsobem kolem osy nebo středu soustavy souřadnic.  
Tuto vlastnost nazýváme \important{sudost} nebo \important{lichost} funkce.  
Určit, zda je funkce sudá či lichá, nám pomáhá lépe pochopit její tvar i chování bez nutnosti vykreslovat celý graf.

\begin{definitionbox}

Funkce \(f\) se nazývá \textbf{sudá}, pokud pro všechna \(x\) z jejího definičního oboru platí:
\[
f(-x) = f(x).
\]
Graf sudé funkce je \textbf{osově souměrný podle osy \(y\)}.
\end{definitionbox}

\begin{examplebox}
\textbf{Příklady sudých funkcí:}
\[
f(x) = x^2, \quad f(x) = \cos x, \quad f(x) = |x|.
\]
Každá z nich má graf, který je při překlápění podle osy \(y\) nezměněný.
\begin{minipage}{0.45\textwidth}
    \begin{center}
        \begin{axisgraph}[grid=none, xmin=-3, xmax=3, ymin=-1, ymax=5, title={Sudá funkce: \(f(x)=x^2\)}]
            \addplot[domain=-3:3, thick, red]{x^2};
        \end{axisgraph}
    \end{center}
    \end{minipage}
    \hfill
    \begin{minipage}{0.45\textwidth}
    \begin{center}
        \begin{axisgraph}[grid=none, xmin=-3, xmax=3, ymin=-3, ymax=3, title={Sudá funkce: \(f(x)=cos(x)\)}]
            \addplot[domain=-3:3, thick, blue]{cos(deg(x))};
        \end{axisgraph}
    \end{center}
    \end{minipage}
\end{examplebox}

\bigskip
\begin{definitionbox}

Funkce \(f\) se nazývá \textbf{lichá}, pokud pro všechna \(x\) z jejího definičního oboru platí:
\[
f(-x) = -f(x).
\]
Graf liché funkce je \textbf{středově souměrný podle počátku soustavy souřadnic}. 
Neboli je středově souměrný podle bodu $[0;0]$
\end{definitionbox}

\begin{examplebox}
\textbf{Příklady lichých funkcí:}
\[
f(x) = x^3, \quad f(x) = \sin x, \quad f(x) = \tan x.
\]
Každý z těchto grafů lze pootočit o \(180^\circ\) kolem počátku a získáme tentýž obraz.
\begin{minipage}{0.45\textwidth}
    \begin{center}
        \begin{axisgraph}[grid=none, xmin=-3, xmax=3, ymin=-3, ymax=3, title={Lichá funkce: \(f(x)=x^3\)}]
            \addplot[domain=-3:3, thick, orange]{x^3};
        \end{axisgraph}
    \end{center}
    \end{minipage}
    \hfill
    \begin{minipage}{0.45\textwidth}
    \begin{center}
        \begin{axisgraph}[grid=none, xmin=-3, xmax=3, ymin=-3, ymax=3, title={Lichá funkce: \(f(x)=sin(x)\)}]
            \addplot[domain=-3:3, thick, purple]{sin(deg(x))};
        \end{axisgraph}
    \end{center}
    \end{minipage}
\end{examplebox}

\begin{notebox}
\textbf{Shrnutí:}
\begin{itemize}
  \item Sudá funkce – „zrcadlí se“ podle osy \(y\), tedy \(f(-x) = f(x)\).
  \item Lichá funkce – „otáčí se“ kolem počátku, tedy \(f(-x) = -f(x)\).
  \item Některé funkce nejsou ani sudé, ani liché (např. \(f(x) = x^2 + x\)).
\end{itemize}
\end{notebox}


% --- P R O S T Á   F U N K C E ---
\subsection*{Prostá funkce}

Každá funkce přiřazuje ke každému vstupu právě jeden výstup,  
ale ne každá to dělá \emph{jedinečně}.  
Může se stát, že dvě různá čísla \(x_1\) a \(x_2\) mají po dosazení stejný výsledek \(f(x_1) = f(x_2)\).  
V takovém případě říkáme o funkci, že \textbf{není prostá}.

Naopak, pokud \textbf{každá výstupní hodnota odpovídá právě jednomé vstupí hodnotě} —  
pak říkáme, že taková funkce je \important{prostá} (neboli \textit{injektivní}).  

Jinými slovy: \textit{prostá funkce se nikdy „nevrací zpět“ a žádná její hodnota se neopakuje.}

\begin{definitionbox}

Funkce \(f\) se nazývá \textbf{prostá}, jestliže pro všechna \(x_1, x_2\) z jejího definičního oboru platí:
\[
f(x_1) \neq f(x_2) \Rightarrow x_1 \neq x_2.
\]
To znamená, že dvě různé vstupní hodnoty nikdy nemohou mít stejný výstup.
\end{definitionbox}
\newpage
Jiným způsobem lze říct, že funkce je prostá právě tehdy, když \textbf{žádná vodorovná přímka neprotne její graf více než v jednom bodě}.  
Tento jednoduchý „vodorovný test“ (angl. \textit{horizontal line test}) je velmi užitečný při určování prostosti z grafu.

\begin{examplebox}
\begin{itemize}
    \item Funkce \(f(x) = x^3\) je prostá, protože pro žádná dvě různá čísla
     \(x_1 \neq x_2\) najdeme vždy dvě různé hodnoty.
    \item Funkce \(f(x) = x^2\) prostá není, protože např. \(f(-2) = f(2) = 4\).  
    Graf má stejnou hodnotu výstupu pro dvě různé vstupní hodnoty.
\end{itemize}

\begin{center}
    \begin{minipage}{0.45\textwidth}
    \centering
        \begin{axisgraph}[
            grid=none,
            xmin=-3, xmax=3,
            ymin=-5, ymax=5,
            title={Prostá funkce $f(x)=x^3$}
        ]
            \addplot[domain=-3:3, thick, blue]{x^3};
            \addplot[dashed, red] coordinates {(-3,4) (3,4)};
        \end{axisgraph}
    \end{minipage}
    \hfill
    \begin{minipage}{0.45\textwidth}
    \centering
        \begin{axisgraph}[
            grid=none,
            xmin=-3, xmax=3,
            ymin=-1, ymax=9,
            title={Neprostá funkce $f(x)=x^2$}
        ]
            \addplot[domain=-3:3, thick, blue]{x^2};
            % vodorovná přímka y=4
            \addplot[dashed, red] coordinates {(-3,4) (3,4)};
        \end{axisgraph}
    \end{minipage}
\end{center}

\end{examplebox}

\begin{notebox}
\textbf{Shrnutí:}
\begin{itemize}
    \item Prostá funkce (\textit{injektivní}) má pro každý výstup právě jeden odpovídající vstup.  
          Žádná hodnota \(y\) se neopakuje.
    \item Neprostá funkce má alespoň dvě různé vstupní hodnoty, které dávají stejný výstup.
    \item Prakticky lze prostost poznat z grafu – \textbf{vodorovná přímka} smí protnout graf \textbf{nejvýše jednou}.
\end{itemize}
\end{notebox}


% --- P E R I O D I C K Á   F U N K C E ---
\subsection*{Periodická funkce}

Zatímco \textbf{prosté funkce} nikdy neopakují stejnou hodnotu dvakrát, tak 
\important{periodické funkce} dělají pravý opak — jejich průběh se po určitém intervalu pravidelně opakuje.  

Takové funkce se chovají „rytmicky“ – jejich graf má tvar, který se po čase znovu vrací.  
Setkáme se s nimi všude tam, kde se děje opakují: u kmitání, vlnění, rotace nebo například u střídání dne a noci.  



\begin{definitionbox}

Funkce \(f\) je \textbf{periodická}, jestliže existuje kladné číslo \(T > 0\), pro které platí:
\[
f(x + T) = f(x) \quad \text{pro všechna } x \text{ z definičního oboru.}
\]
Číslo \(T\) zde nazýváme \textbf{perioda funkce}.
\end{definitionbox}

To znamená, že posuneme-li graf funkce o délku \(T\) ve směru osy \(x\), dostaneme \emph{tentýž} graf.  
Každý bod se tedy po určité „vzdálenosti“ znovu zopakuje.  
Periodické funkce se v matematice i v přírodě vyskytují velmi často — popisují pohyby, kmitání, oscilace, střídání veličin či signálů.

\begin{examplebox}
\textbf{Příklady periodických funkcí:}
\begin{itemize}
    \item Funkce \(f(x) = \sin x\) má periodu \(T = 2\pi\).
    \item Funkce \(f(x) = \cos x\) má periodu \(T = 2\pi\).
    \item Funkce \(f(x) = \tan x\) má periodu \(T = \pi\).
    \item Funkce \(f(x) = | \sin x |\) má periodu \(T = \pi\).
\end{itemize}
\begin{center}
\begin{axisgraph}[
    width=12cm,     % šířka grafu
    height=6cm,
    grid=minor,
    xmin=-3.5, xmax=13,
    ymin=-1.5, ymax=1.5,
    title={Periodická funkce $f(x) = \cos x$}
]
    \addplot[domain=-3.5:13, thick, blue, samples=200]{cos(deg(x))};
    \addplot[domain=0:2*pi, thick, red, samples=200]{cos(deg(x))};
    
    % pomocné značky pro periodu
    \draw[dashed, red] (2*pi,-1.5) -- (2*pi,1.5);
    \node[red, below] at (2*pi, -1.5) {$T = 2\pi$};
\end{axisgraph}
\end{center}
Jak můžeme vidět na grafu $f(x)=cos(x)$, tak červená část grafu se nám poté stále opakuje.
\end{examplebox}



\begin{notebox}
\textbf{Shrnutí:}
\begin{itemize}
    \item Periodická funkce se po určité délce \(T\) opakuje: \(f(x+T) = f(x)\).
    \item Nejmenší takové kladné \(T\) se nazývá \textbf{základní perioda}.
    \item Typickými příklady jsou goniometrické funkce (sinus, kosinus, tangens).
    \item V praxi se periodické funkce používají při popisu dějů, které se pravidelně opakují —  
          například střídavé napětí, pohyb kyvadla nebo zvukové vlny.
\end{itemize}
\end{notebox}


% --- O M E Z E N O S T   F U N K C E ---
\subsection*{Omezenost funkce}

V matematice často zkoumáme, \textbf{jak vysoko nebo nízko může graf funkce „dosáhnout“}.  
Některé funkce se rozbíhají do nekonečna (např. \(f(x) = x^2\) stále roste, když zvětšujeme \(x\)),  
jiné naopak zůstávají „v nějakém pásmu“ a nikdy nepřekročí určitou hodnotu.  
Takové funkce nazýváme \important{omezené}. A máme funkce buď \important{shora omezené}, nebo \important{zdola omezené}.

\begin{definitionbox}

Funkce je \textbf{omezená shora}, jestliže existuje číslo \(A\), pro které platí:
    \[
    f(x) \leq A \quad \text{pro všechna } x \text{ z definičního oboru.}
    \]
\emph{Říkáme, že \(A\) je \textbf{horní závora} funkce.}

\smallskip
Funkce je \textbf{omezená zdola}, jestliže existuje číslo \(a\), pro které platí:
    \[
    f(x) \geq a \quad \text{pro všechna } x \text{ z definičního oboru.}
    \]
\emph{Číslo \(a\) pak nazýváme \textbf{dolní závora} funkce.}
    
Jestliže existují obě závory \(a\) i \(A\), pak říkáme, že funkce je \textbf{omezená}.

\end{definitionbox}

Intuitivně bychom mohli říct, že graf funkce tzv. \emph{„zůstává mezi dvěma vodorovnými přímkami“} –  
nikdy nepřekročí horní hranici a nikdy neklesne pod dolní hranici.  
Omezenost je tedy vlastnost, která říká, zda funkce „utíká“ do nekonečna, nebo se drží v určité oblasti.


\begin{examplebox}
Funkce \(f(x) = \sin x\) je omezená shora i zdola, protože platí \(-1 \leq \sin x \leq 1\) pro všechna \(x \in \mathbb{R}\).
\begin{center}
    \begin{axisgraph}[grid=minor, xmin=-6.5, xmax=6.5, ymin=-2, ymax=2, title={Graf funkce: \(f(x)=\sin x\)}]
        \addplot[domain=-6.5:6.5, thick, magenta]{sin(deg(x))};
        % vodorovná přímka y=4
            \addplot[dashed, red] coordinates {(-6.5,-1.5) (6.5,-1.5)};
            \addplot[dashed, red] coordinates {(-6.5,1.5) (6.5,1.5)};
    \end{axisgraph}
\end{center}
Naopak graf funkce $f(x) = x^3$ omezený není, protože můžeme najít pro všechna $x$ z této funkce, vždy vyšší hodnotu.
\begin{center}
    \begin{axisgraph}[grid=minor, xmin=-3, xmax=3, ymin=-4, ymax=4, title={Graf funkce: \(f(x)=x^3\)}]
        \addplot[domain=-5:5, thick, cyan]{x^3};
    \end{axisgraph}
\end{center}
\end{examplebox}

\begin{notebox}
\textbf{Shrnutí:}
\begin{itemize}
    \item Funkce omezená shora: má „strop“ – existuje maximální dosažitelná hodnota.
    \item Funkce omezená zdola: má „podlahu“ – existuje minimální dosažitelná hodnota.
    \item Funkce omezená: zůstává mezi dvěma hodnotami, tj. \(m \leq f(x) \leq M\).
\end{itemize}
Všimněte si, že vůbec nemusíme najít taková $A$, $a$, které jsou „nejbližší“ grafu, 
stačí přímku umístit libovolně pod/nad graf.
\end{notebox}


% --- M A X I M U M   A   M I N I M U M ---
\subsection*{Maximum a minimum}

V průběhu funkce často hledáme body, které dosahují \textbf{nejvyšší} nebo \textbf{nejnižší} hodnoty.  
Tyto body nazýváme \important{extrémy funkce}.  
Extrémy jsou velmi důležité, protože ukazují, kde se funkce „otáčí“ – přestává růst a začíná klesat, nebo naopak.  
V praxi mohou představovat například největší zisk, nejmenší náklady, nejvyšší rychlost či nejnižší teplotu.

\begin{definitionbox}

    Funkce má v bodě \(x_0\) \textbf{lokální maximum}, jestliže platí
    \[
    f(x_0) \geq f(x)
    \quad \text{pro všechna } x \text{ z okolí bodu } x_0.
    \]
    V takovém bodě má funkce „vrchol“ ve svém okolí.

    \medskip
    Funkce má v bodě \(x_0\) \textbf{lokální minimum}, jestliže platí
    \[
    f(x_0) \leq f(x)
    \quad \text{pro všechna } x \text{ z okolí bodu } x_0.
    \]
    Takový bod představuje „údolí“ funkce v jejím okolí.

    Jestliže funkce dosáhne největší nebo nejmenší hodnoty na \emph{celém definičním oboru},  
    nazýváme tyto body \textbf{globální maximum} a \textbf{globální minimum}.
\end{definitionbox}

\begin{center}
\begin{minipage}{0.49\textwidth}
    \small
    Graf vpravo zobrazuje funkci \(f(x) = x^2 - 1\),  
    která má tvar \textbf{otevřené paraboly směrem vzhůru}.  
    V bodě \(x = 0\) dosahuje \textbf{své nejnižší hodnoty} –  
    to je \important{globální minimum} \(f(0) = -1\).  

    Funkce nemá žádné maximum,  
    protože s rostoucím \(|x|\) její hodnota neomezeně roste.  
    Tento tvar je typický pro kvadratické funkce s kladným koeficientem u \(x^2\).
\end{minipage}
\hfill
\begin{minipage}{0.5\textwidth}
    \begin{axisgraph}[
        grid=none,
        xmin=-3, xmax=3,
        ymin=-2, ymax=4,
        width=\textwidth,
        title={Parabola \(f(x) = x^2 - 1\)}
    ]
        \addplot[domain=-3:3, thick, blue]{x^2 - 1};
        % zvýraznění minima
        \addplot[mark=*] coordinates {(0,-1)};
        \node[below right] at (axis cs:0,-1) {\small $[0, -1]$};
    \end{axisgraph}
\end{minipage}%
\end{center}

\begin{notebox}
\textbf{Shrnutí:}
\begin{itemize}
    \item \textbf{Lokální extrémy} se týkají chování funkce v blízkosti určitého bodu.
    \item \textbf{Globální extrémy} se týkají celé funkce – určují její absolutně nejvyšší a nejnižší hodnotu.
    \item Parabola \(f(x) = x^2 - 1\) má jedno \textbf{globální minimum} v bodě \([0, -1]\)  
          a žádné maximum, protože roste do nekonečna.
\end{itemize}
\end{notebox}


% --- K O N V E X N O S T   A   K O N K Á V N O S T ---
\subsection*{Konvexnost a konkávnost}

Kromě toho, zda funkce \textbf{roste} nebo \textbf{klesá}, nás často zajímá také to, \emph{jakým způsobem} se její graf ohýbá.  
Jinými slovy – zda se křivka „usmívá“ nebo „mračí“.  
Tento tvar grafu popisujeme pojmy \important{konvexnost} a \important{konkávnost}.

\begin{definitionbox}

    Funkce \(f(x)\) je \textbf{konvexní}, pokud leží celý její graf \textbf{pod všemi přímkami}, které spojují libovolné dva body grafu.

    Funkce \(f(x)\) je \textbf{konkávní}, pokud leží celý její graf \textbf{nad těmito přímkami}.
\end{definitionbox}

\textbf{Intuitivně:}

\begin{itemize}
  \item \textbf{Konvexní funkce} má tvar \emph{misky} – otevírá se vzhůru, „usmívá se“.  
  Každá tečna (nebo spojovací přímka) leží nad grafem.  
  Typický příklad: \(f(x) = x^2\).

  \item \textbf{Konkávní funkce} má tvar \emph{kopce} – otevírá se dolů, „mračí se“.  
  Každá tečna leží pod grafem.  
  Typický příklad: \(f(x) = -x^2\).
\end{itemize}

\begin{center}
\begin{axisgraph}[
    width=8cm,     % šířka grafu
    height=8cm,
    grid=none,
    xmin=-2, xmax=3,
    ymin=-1.5, ymax=3,
    title={Konvexní a konkávní část grafu: $f(x)=x^3 - 2\cdot x^2 - x + 2$}
]
    % konkávní část (vlevo)
    \addplot[domain=-3:0.6, thick, orange]{x^3 - 2*x^2 - x + 2};
    % konvexní část (vpravo)
    \addplot[domain=0.6:3, thick, blue]{x^3 - 2*x^2 - x + 2};
    % úsečky k bodu
    \addplot[dashed, black] coordinates {(0.6,0) (0.6,0.94)};
    \addplot[dashed, black] coordinates {(0,0.9) (0.6,0.9)};
    % vyznačení inflexního bodu
    \addplot[mark=*, mark size=2.5pt, color=red] coordinates {(0.6,0.9)};
    \node[red, above right] at (axis cs:0.6,1) {\small inflexní bod};
    \node[black, above right] at (axis cs:0.6,0.7) {\scriptsize $[0,6;0,9]$};
    % popisky
    \node[orange!90!black] at (axis cs:-1.5,1.2) {\small konkávní};
    \node[blue!80!black] at (axis cs:1.5,0.5) {\small konvexní};
    \node[black] at (axis cs:0.6,-0.15) {\small 0,6};
    \node[black] at (axis cs:-0.3,0.8) {\small 0,9};
\end{axisgraph}
\end{center}

Funkce $f(x)=x^3 - 2\cdot x^2 - x + 2$ je v levé části konkávní, v pravé konvexní.  
V bodě \([0,6;0,9]\) se křivka „láme“ – tam má \textbf{inflexní bod}.

\begin{notebox}
\textbf{Shrnutí:}
\begin{itemize}
  \item Konvexní funkce má tvar „misky“ – \textit{směrem nahoru}, tečny leží nad ní.  
  \item Konkávní funkce má tvar „kopce“ – \textit{směrem dolů}, tečny leží pod ní.  
  \item V bodě, kde se mění konvexnost na konkávnost (nebo naopak), má funkce \textbf{bod inflexe} – tam se křivka „přelamuje“.
\end{itemize}
\end{notebox}





% --- M O N O T Ó N N O S T   F U N K C E ---
\subsection*{Monotónnost}

\important{Monotónnost} popisuje, \textbf{jak se mění hodnota funkce, když se zvětšuje argument \(x\)}.  
Jinými slovy — zkoumáme, zda funkce roste, klesá, nebo zůstává stále stejná.  
Podle toho rozlišujeme tři základní typy monotónnosti:

\begin{definitionbox}
Funkce \(f\) je:
\begin{itemize}
    \item \textbf{rostoucí}, jestliže pro všechna \(x_1 < x_2\) platí \(f(x_1) < f(x_2)\);  
          s rostoucím \(x\) se tedy zvyšuje i hodnota funkce \(f(x)\);
    \item \textbf{klesající}, jestliže pro všechna \(x_1 < x_2\) platí \(f(x_1) > f(x_2)\);  
          čím větší je \(x\), tím menší je \(f(x)\);
    \item \textbf{konstantní}, jestliže \(f(x_1) = f(x_2)\) pro všechna \(x_1, x_2\) z definičního oboru;  
          hodnota funkce se nemění.
\end{itemize}
\end{definitionbox}

Monotónnost je velmi důležitá vlastnost — pomáhá nám poznat chování funkce, určovat její extrémy, a lépe chápat tvar grafu.  
Na rostoucím intervalu se graf „zvedá“, na klesajícím se „svažuje“ a u konstantní funkce tvoří vodorovnou přímku.

\begin{center}
\begin{axisgraph}[
    grid=none,
    xmin=-3, xmax=5,
    ymin=-2, ymax=3,
    width=12cm,
    height=7cm,
    title={Různé úseky monotónnosti funkce}
]
    % Klesající část
    \addplot[domain=-2:-1, very thick, red]{-2*x - 3};      % klesání z (−2,3) do (−1,−1)
    % Konstantní část
    \addplot[domain=-1:2, thick, gray]{-1};             % konstantní úsek
    % Rostoucí část
    \addplot[domain=2:4, thick, blue]{1.5*x - 4};       % rostoucí úsek z (2,−1) do (4,2)

    % body přechodu
    \addplot[mark=*] coordinates {(-2,1) (-1,-1) (2,-1) (4,2)};
    
    % popisky oblastí
    \node[red!90!black] at (axis cs:-1.5,1.5) {\small klesající};
    \node[gray!70!black]   at (axis cs:0.5,-0.7) {\small konstantní};
    \node[blue!80!black]   at (axis cs:3,0.2) {\small rostoucí};

\end{axisgraph}
\end{center}

O grafu výše můžeme říct, že na intervalu $(-2;-1)$ \textbf{klesá}, $(-1;2)$ je \textbf{konstantní} a $(2;4)$ \textbf{roste}.


% ===========================================
% Z Á K L A D N Í   E L E M E N T Á R N Í   F U N K C E
% ===========================================

\section{Základní elementární funkce}

V této části se seznámíme s nejběžnějšími typy funkcí, se kterými se budeme v dalších kapitolách opakovaně setkávat.  
U každé z nich si uvedeme základní tvar předpisu, vzhled grafu a stručnou charakteristiku jejích vlastností.

% --- L I N E Á R N Í   F U N K C E ---
\subsection*{Lineární funkce}

\begin{definitionbox}
\textbf{Lineární funkce} má předpis
\[
f(x) = a x + b,
\]
kde \(a\) je \textbf{směrnice přímky} a \(b\) je \textbf{posunutí} (průsečík s osou \(y\)).
\end{definitionbox}

Lineární funkce představuje přímku v souřadnicové rovině.  
Její tvar závisí na hodnotě směrnice \(a\):  
pokud \(a > 0\), funkce roste; pokud \(a < 0\), funkce klesá; a pro \(a = 0\) je konstantní.  
Hodnota \(b\) posouvá přímku nahoru nebo dolů.

\begin{center}
\begin{axisgraph}[grid=none, xmin=-3, xmax=3, ymin=-3, ymax=3, title={Příklady lineárních funkcí}]
    \addplot[domain=-3:3, thick, blue]{x};           % rostoucí
    \addplot[domain=-3:3, thick, orange]{-0.5*x+1};  % klesající
    \addplot[domain=-3:3, thick, gray]{1};           % konstantní

    \node[blue] at (axis cs:2.3,1.5) {\small rostoucí};
    \node[orange] at (axis cs:-1.5,2.3) {\small klesající};
    \node[gray] at (axis cs:-2,0.7) {\small konstantní};
\end{axisgraph}
\end{center}

\begin{notebox}
\textbf{Shrnutí:}  
Lineární funkce je nejjednodušším typem funkce – vyjadřuje přímou úměrnost nebo její posunutou variantu.  
Grafem je přímka a její chování zcela určuje směrnice \(a\).
\end{notebox}


% --- K V A D R A T I C K Á   F U N K C E ---
\subsection*{Kvadratická funkce}

\begin{definitionbox}
\textbf{Kvadratická funkce} má tvar
\[
f(x) = a x^2 + b x + c, \quad a \neq 0.
\]
\end{definitionbox}

Grafem kvadratické funkce je \textbf{parabola}.  
Její osa symetrie má směr rovnoběžný s osou \(y\).  
Pokud \(a > 0\), parabola je otevřená nahoru, pokud \(a < 0\), je otevřená dolů.  
Vrchol paraboly se nachází v bodě
\[
[x_v, y_v] = \left[-\frac{b}{2a}, f\!\left(-\frac{b}{2a}\right)\right].
\]

\begin{center}
\begin{axisgraph}[grid=none, xmin=-4, xmax=4, ymin=-2, ymax=6, title={Příklady kvadratických funkcí}]
    \addplot[domain=-3:3, thick, blue]{x^2};
    \addplot[domain=-5:5, thick, orange]{-0.5*x^2+4};
\end{axisgraph}
\end{center}


% --- M O C N I N N Á   F U N K C E ---
\subsection*{Mocninná funkce}

\begin{definitionbox}
\textbf{Mocninná funkce} má tvar
\[
f(x) = x^n,
\]
kde \(n\) je reálné číslo.
\end{definitionbox}

Vlastnosti závisí na hodnotě exponentu \(n\):
\begin{itemize}
  \item pro \(n\) sudé je funkce \textbf{sudá} a má tvar „misky“ (např. \(x^2, x^4\)),
  \item pro \(n\) liché je funkce \textbf{lichá} a má tvar „eska“ (např. \(x^3, x^5\)).
\end{itemize}

\begin{center}
\begin{axisgraph}[grid=none, xmin=-3, xmax=3, ymin=-3, ymax=3, title={Příklady mocninných funkcí}]
    \addplot[domain=-3:3, thick, blue]{x^2};
    \addplot[domain=-3:3, thick, orange]{x^3};
\end{axisgraph}
\end{center}


% --- L O M E N Á   F U N K C E ---
\subsection*{Lomená funkce}

\begin{definitionbox}
\textbf{Lomená funkce} má tvar
\[
f(x) = \frac{a}{x} \quad \text{nebo obecně} \quad f(x) = \frac{a x + b}{c x + d}.
\]
\end{definitionbox}

Její graf se nazývá \textbf{hyperbola}.  
Funkce není definována pro hodnoty, které nulují jmenovatel.  
Hyperbola má dvě \textbf{asymptoty} – přímky, ke kterým se přibližuje, ale nikdy je neprotne.

\begin{center}
\begin{minipage}{0.45\textwidth}
\centering
\begin{axisgraph}[
    grid=none,
    xmin=-5, xmax=5,
    ymin=-5, ymax=5,
    title={Základní tvar $f(x)=\frac{1}{x}$}
]
    % Funkce 1/x
    \addplot[domain=-5:-0.2, thick, blue]{1/x};
    \addplot[domain=0.2:5, thick, blue]{1/x};
    % Asymptoty
    \addplot[dashed, gray]{0}; % y=0
    \addplot[dashed, gray] coordinates {(0,-5) (0,5)}; % x=0
\end{axisgraph}
\end{minipage}
\hspace{1cm}
\begin{minipage}{0.45\textwidth}
\centering
\begin{axisgraph}[
    grid=none,
    xmin=-5, xmax=5,
    ymin=-5, ymax=5,
    title={Posunutý tvar $f(x)=\frac{x+3}{x-2}$}
]
    % Funkce (x+3)/(x-2)
    \addplot[domain=-5:1.8, thick, orange]{(x+3)/(x-2)};
    \addplot[domain=2.2:5, thick, orange]{(x+3)/(x-2)};
    % Asymptoty
    \addplot[dashed, gray]{1}; % vodorovná asymptota y=1
    \addplot[dashed, gray] coordinates {(2,-5) (2,5)}; % svislá asymptota x=2
    % Popisky asymptot
    \node[gray] at (axis cs:2.3,4.2) {\small $x=2$};
    \node[gray] at (axis cs:-4.2,1.2) {\small $y=1$};
\end{axisgraph}
\end{minipage}
\end{center}



% --- A B S O L U T N Í   H O D N O T A ---
\subsection*{Funkce s absolutní hodnotou}

\begin{definitionbox}
\textbf{Funkce s absolutní hodnotou} má tvar
\[
f(x) = |x|.
\]
\end{definitionbox}

Grafem je „zlomená přímka“ ve tvaru písmene \(V\).  
Funkce je vždy nezáporná, roste pro \(x>0\) a klesá pro \(x<0\).  
Je \textbf{sudá} a má minimum v bodě \((0,0)\).

\begin{center}
\begin{axisgraph}[grid=none, xmin=-3, xmax=3, ymin=-1, ymax=3, title={Funkce s absolutní hodnotou $f(x)=|x|$}]
    \addplot[domain=-3:3, thick, blue]{abs(x)};
\end{axisgraph}
\end{center}


% -----------------------------------------------------------------
% ---------------------- P Ř Í K L A D Y --------------------------
% -----------------------------------------------------------------

\exercisepart
\setcounter{section}{1}


% =========================
% ÚLOHA 1
% =========================
\begin{exercise}[Zapiš funkční závislosti I]
    Zapište jako funkce (s vhodně zvolenou proměnnou) následující závislosti; dbejte na to, aby byl zápis jasný, jednoznačný a aby byla proměnná výslovně popsána:
    \begin{tasks}(2)
        \task obsah čtverce na délce jeho strany,
        \task obvod kruhu na jeho poloměru,
        \task obsah kruhu na jeho průměru,
        \task objem krychle na délce její strany,
        \task obsah obdélníku v závislosti na délce kratší strany, přičemž delší strana je \(2{,}5\!\times\) delší než strana kratší,
        \task velikost tělesové úhlopříčky krychle na délce její strany.
    \end{tasks}
\end{exercise}


% =========================
% ÚLOHA 2
% =========================
\begin{exercise}[Zapiš funkční závislosti II]
    Zapište jako funkce (s vhodně zvolenou proměnnou) následující závislosti; dbejte na to, aby byl zápis jasný, jednoznačný a aby byla proměnná výslovně popsána:
    \begin{tasks}(1)
        \task obvodu rovnoramenného pravoúhlého trojúhelníku na délce $a$ jeho odvěsny, 
        \task obvodu rovnoramenného pravoúhlého trojúhelníku na délce $c$ jeho přepony.
    \end{tasks}
\end{exercise}


% =========================
% ÚLOHA 3
% =========================
\begin{exercise}[Kvádr se čtvercovou podstavou]
    Uvažujte kvádr se čtvercovou podstavou; délka podstavné hrany je \(b\) a délka boční (svislé) hrany je rovna polovině podstavné hrany. Zapište funkce udávající závislost následujících veličin na \(b\) a u každé stručně vysvětlete postup:
    \begin{tasks}(2)
        \task součet délek hran kvádru na \(b\),
        \task délku tělesové úhlopříčky na \(b\),
        \task povrch kvádru na \(b\),
        \task objem kvádru na \(b\).
    \end{tasks}
\end{exercise}


% =========================
% ÚLOHA 4
% =========================
\begin{exercise}[Rozpoznání, zda je daný předpis funkcí]
    Rozhodněte, zda následující předpisy určují funkci na uvedené množině nezávisle proměnné \(x\).  

    \begin{tasks}(2)
        \task \(f_{1}:~ y = 2x - 1,\quad x \in \mathbb{R}\)
        \task \(f_{2}:~ y = \sqrt{x + 4},\quad x \in \langle -4; \infty )\)
        \task \(f_{3}:~ y^{2} = x,\quad \quad \quad x \in \mathbb{R}\)
        \task \(f_{4}:~ y = x^{2},\quad \quad \quad x \in \mathbb{R}\)
    \end{tasks}

\end{exercise}


% =========================
% ÚLOHA 5
% =========================
\begin{exercise}[Definiční obory elementárních funkcí]
    Určete definiční obor každé z následujících funkcí:
    \begin{tasks}(2)
        \task \(f_{1}(x)=x^{2}+3x-1\)
        \task \(f_{2}(x)=\dfrac{x+1}{x-3}\)
        \task \(f_{3}(x)=\dfrac{x}{13x^{2}+10x-3}\)
        \task \(f_{4}(x)=\dfrac{x^{2}+1}{x^{2}+x+1}\)
        \task \(f_{5}(x)=\dfrac{1}{\,|x+3|-4\,}\)
        \task \(f_{6}(x)=\sqrt{x+2}\)
    \end{tasks}
\end{exercise}


% =========================
% ÚLOHA 6
% =========================
\begin{exercise}[Výpočet funkčních hodnot I]
    Pro funkci \(y=2x^{2}-3x\), kde \(x\in\mathbb{R}\), vypočtěte přesné (nebo vhodně zjednodušené) hodnoty v bodech
    \(-4\), \(0{,}3\), \(\sqrt{7}\) a \(12\).
\end{exercise}


% =========================
% ÚLOHA 7
% =========================
\begin{exercise}[Výpočet funkčních hodnot II]
    Je dána funkce \(f\colon y=\dfrac{5-x}{x+2}\) s omezením \(x\in\langle -1;5\rangle\). Vypočtěte hodnoty
    \(f(-1)\), \(f(0{,}6)\), \(f(2)\) a \(f(4{,}2)\). Pokud výraz není pro zadaný bod definován, uveďte důvod.
\end{exercise}


% =========================
% ÚLOHA 8
% =========================
\begin{exercise}[Výpočet funkčních hodnot]
    Je dána funkce $g:y=\frac{2}{x^2 - 1}$.
    \begin{tasks}
        \task Určete definiční obor funkce $g$ pomocí sjednocení intervalů.
        \task Určete u(0); u(3); u(-7) u(1,5).
        \task Která z čísel $-4;0;1,8$ patří do oboru hodnot funkce g.
    \end{tasks}
\end{exercise}


% =========================
% ÚLOHA 9
% =========================
\begin{exercise}[Hodnota funkce]
    Vypočítejte hodntotu funkce $f:y=\frac{x+4}{x-1}$ v bodech $-1;3;\sqrt{2}$.

\end{exercise}


% =========================
% ÚLOHA 10
% =========================
\begin{exercise}[Určení hodnot funkce z grafu]
    Mějme graf funkce \(t\), jejíž definiční obor je interval \(\langle -4;4\rangle\).
    \begin{tasks}(1)
        \task Určete hodnoty funkce \(t\) alespoň v osmi různých bodech.
        \task Zapište obor hodnot funkce \(t\) pomocí intervalu.
        \task Zapište všechna \(x \in D(t)\), pro která je $ t(x)=-2,\; t(x)=0,\; t(x)=1,\; t(x)=2,\; t(x)\geq 0.$
    \end{tasks}
    \begin{center}
        \begin{axisgraph}[
            xmin=-4, xmax=4,
            ymin=-2, ymax=2,
            width=12cm,
            height=6cm,
            axis lines=middle,
            grid=none,
            xtick={-4,-3,-2,-1,0,1,2,3,4},
            ytick={-2,-1,0,1,2},
            minor tick num=1
        ]
        % --- místo pro graf funkce t(x) ---
        % \addplot[thick, smooth, domain=-4:4, samples=100] { ... };
        \node at (rel axis cs:0.92,0.9) {$t$};
        % body přechodu
        \addplot[very thick,mark=*] coordinates {(-4,1) (-3,1) (-2,2) (0,-2) (2,2) (3,1) (4,1)};
        \end{axisgraph}
    \end{center}
\end{exercise}


% =========================
% ÚLOHA 11
% =========================
\begin{exercise}[Základní práce s definičním oborem a oborem hodnot]
    Mějme funkci \(h:\ y=\dfrac{2}{x^{2}+1}\).
    \begin{tasks}(1)
        \task Určete její definiční obor.
        \task Vypočtěte hodnotu funkce \(h\) v bodě \(x=3\).
        \task Rozhodněte, zda \(4 \in H(h)\), tedy zda číslo \(4\) náleží do oboru hodnot této funkce; své rozhodnutí krátce zdůvodněte.
    \end{tasks}
\end{exercise}


% =========================
% ÚLOHA 12
% =========================
\begin{exercise}[Definiční obor složené odmocninové funkce s omezením]
    Je dána funkce \(f:\ y=\sqrt{x-3}+\sqrt{\,5-x\,}\) s doplňkovým omezením nezávisle proměnné \(x\in\langle -6;9\rangle\).
    Určete její definiční obor tak, že zkombinujete podmínky plynoucí z pododmocninových výrazů s uvedeným intervalovým omezením.
\end{exercise}


% =========================
% ÚLOHA 13
% =========================
\begin{exercise}[Výpočet funkčních hodnot I]
    Je dána funkce \(f:\ y = x^{3} - x,\quad x \in \mathbb{R}\).  
    Vypočítejte hodnoty funkce pro zadané argumenty:
 $f(-2), f\!\left(\tfrac{1}{3}\right), f(2\sqrt{3}).$
\end{exercise}


% =========================
% ÚLOHA 14
% =========================
\begin{exercise}[Výpočet funkčních hodnot II]
    Je dána funkce \(g(x)=3x^{2}+1\).  
    Určete funkční hodnoty pro zadané argumenty: $g(1), g(a)+g(2), g(a+2), g(b^{2}), [\,g(b)\,]^{2}$.
\end{exercise}


% =========================
% ÚLOHA 15
% =========================
\begin{exercise}[Existence řešení pro dané hodnoty funkce]
    Je dána funkce \(f(x)=2x-9\).  
    Rozhodněte, zda existuje takové \(x\in\mathbb{R}\), aby platilo:
    \begin{tasks}(3)
        \task \(f(x)=5\)
        \task \(f(x)=-100\)
        \task \(f(x)=-\dfrac{2}{3}\)
    \end{tasks}
\end{exercise}



% =========================
% ÚLOHA 16
% =========================
\begin{exercise}[Rovnost funkcí a role definičního oboru]
    Rozhodněte, zda jsou si rovny funkce \(f:\ y=1\) a \(g:\ y=\dfrac{x}{x}\).
    Uveďte jak porovnání hodnot, tak i diskusi o jejich definičních oborech, protože rovnost funkcí vyžaduje shodu nejen obrazů, ale také shodu předpisu na totožném definičním oboru.
\end{exercise}


% =========================
% ÚLOHA 17
% =========================
\begin{exercise}[Práce s grafy tří funkcí]
    Na obrázku jsou grafy funkcí \(g_{1}\), \(g_{2}\) a \(g_{3}\).
        \begin{center}
    % --- tři samostatné grafy vedle sebe ---
    \begin{minipage}[t]{0.32\textwidth}
        \vspace{0pt}
        \begin{axisgraph}[
            xmin=-2.5, xmax=3.5,
            ymin=-2, ymax=3,
            width=\linewidth,
            height=6cm,
            axis lines=middle,
            grid=none,
            xtick={-4,-2,0,3},
            ytick={1,2},
            minor tick num=1
        ]
        % --- sem vlož křivku g1, např.: \addplot[thick] { ... }; ---
        \node at (rel axis cs:0.92,0.9) {$g_{1}$};
        % --- úsečka: [-2,2] až (0,2) s otevřeným koncem v (0,2) ---
        \addplot[thick] coordinates {(-2,2) (-0.18,2)};
        \addplot[only marks, mark=*, mark size=2pt] coordinates {(-2,2)}; % uzavřený bod
        \addplot[only marks, mark=o, mark size=2pt] coordinates {(0,2)};  % otevřený bod
        % --- úsečka: [0,0] až [3,1] (oba konce zahrnuté) ---
        \addplot[thick] coordinates {(0,0) (3,1)};
        \addplot[only marks, mark=*, mark size=2pt] coordinates {(0,0) (3,1)};
        \end{axisgraph}
    \end{minipage}\hfill
    \begin{minipage}[t]{0.32\textwidth}
        \vspace{0pt}
        \begin{axisgraph}[
            xmin=-3, xmax=3,
            ymin=-2.1, ymax=3,
            width=\linewidth,
            height=6cm,
            axis lines=middle,
            grid=none,
            xtick={-1},
            ytick={2},
            minor tick num=1
        ]
        % --- sem vlož křivku g2 ---
        \node at (rel axis cs:0.92,0.9) {$g_{2}$};
        % --- část kružnice se středem (-2,0) ---
        \addplot[domain=-2:0, samples=100, thick] ({x},{-sqrt(4-(x+2)^2)});
                
        % --- část kružnice se středem (2,0) ---
        \addplot[domain=0:2, samples=100, thick] ({x},{-sqrt(4-(x-2)^2)+2});
                
        % --- body ---
        \addplot[only marks, mark=*, mark size=2pt] coordinates {(0,2)}; % uzavřený bod nahoře
        \addplot[only marks, mark=o, mark size=2pt] coordinates {(0,0)}; % otevřený bod dole
        \end{axisgraph}
    \end{minipage}\hfill
    \begin{minipage}[t]{0.32\textwidth}
        \vspace{0pt}
        \begin{axisgraph}[
            xmin=-3, xmax=3,
            ymin=-2, ymax=3,
            width=\linewidth,
            height=6cm,
            axis lines=middle,
            grid=none,
            xtick={2,4},
            ytick={1/2},
            minor tick num=1
        ]
        % --- sem vlož křivku g3 ---
        \node at (rel axis cs:0.92,0.9) {$g_{3}$};
        \addplot[dashed, gray] coordinates {(-1,-5) (-1,5)};
        % --- část kružnice se středem (2,0) ---
        \addplot[domain=-3:0, samples=100, thick] ({x},{sqrt(4-(x+3.2)^2)-2.1});

        \addplot[thick] coordinates {(0,1/2) (2,1/2)};
        \addplot[only marks, mark=*, mark size=2pt] coordinates {(0,1/2) (2,1/2)};
        \end{axisgraph}
    \end{minipage}
    % --- konec trojice grafů ---
    \end{center}
    \begin{tasks}(1)
        \task Určete definiční obor a obor hodnot pro funkce \(g_{1}\), \(g_{2}\) a \(g_{3}\).
        \task Určete čísla \(x_{1}, x_{2}, x_{3}\) z definičního oboru funkce \(g_{1}\) tak, aby platilo: \\
            $g_{1}(x_{1})=0{,}5,\qquad g_{1}(x_{2})=2,\qquad g_{1}(x_{3})=-3.$
    \end{tasks}
\end{exercise}


% =========================
% ÚLOHA 18
% =========================
\begin{exercise}[Graf závislosti rychlosti na vzdálenosti od Země]
    Má-li umělá družice kroužit kolem Země a neklesat, musí dosáhnout jisté minimální rychlosti, tzv.\ první kosmické rychlosti. V následující tabulce jsou uvedeny dvojice zachycující závislost číselné hodnoty rychlosti \(v\) na číselné hodnotě vzdálenosti \(h\) družice od Země (vzdálenost je vyjádřena v kilometrech a rychlost v~\(\mathrm{km\,s^{-1}}\)):

    \medskip
    \centering
    \begin{tabular}{c|cccccccc}
        \(h\) [km]       & 50 & 100 & 200 & 300 & 400 & 500 & 600 & 700 \\
        \hline
        \(v\) [km\,s\(^{-1}\)] & 7{,}6 & 7{,}4 & 6{,}9 & 6{,}5 & 5{,}2 & 5{,}9 & 5{,}7 & 5{,}4
    \end{tabular}
    \medskip

    V pravoúhlé soustavě souřadnic pečlivě znázorněte všechny uspořádané dvojice čísel uvedené v tabulce tak, že na vodorovnou osu nanesete hodnoty \(h\) a na svislou osu hodnoty \(v\), přičemž zvolíte rozumné měřítko, označíte obě osy s jednotkami a body vyznačíte jednoznačně tak, aby bylo zřejmé, jak se rychlost mění se vzrůstající vzdáleností.
\end{exercise}


% =========================
% ÚLOHA 19
% =========================
\begin{exercise}[Modelová data: přírůstky objemu vody při povodni a kumulace v čase]
    Uvažujme zjednodušený model průběhu povodně, v němž sledujeme \emph{přírůstek objemu} vody v přehradní nádrži v jednotlivých hodinách; v tabulce jsou uvedeny modelové (didakticky zjednodušené) hodnoty přírůstku \(\Delta V\) v~milionech metrů krychlových \([\mathrm{mil.\ m^3}]\) během postupných hodin \([h]\) od začátku srážkové epizody:

    \medskip
    \centering
    \begin{tabular}{c|ccccccccccc}
        \(t\) [h]          & 0 & 1 & 2 & 3 & 4 & 5 & 6 & 7 & 8 & 9 & 10\\
        \hline
        \(\Delta V\) [mil.\,m\(^3\)] & 0{,}10 & 0{,}15 & 0{,}22 & 0{,}26 & 0{,}31 & 0{,}28 & 0{,}24 & 0{,}20 & 0{,}17 & 0{,}14 & 0{,}11
    \end{tabular}
    \medskip

    \begin{tasks}
        \task Sestavte funkci \(S\), která každé celé hodině \(t=0,1,2,\dots,12\) přiřadí \emph{celkový} (kumulativní) přírůstek objemu od začátku epizody.
        \task V pravoúhlé soustavě souřadnic znázorněte uspořádané dvojice \((t, V)\) a graf popište (osy, jednotky, rozumné měřítko).
        \task Diskutujte, zda je \(S\) na dané nosné množině rostoucí funkce a vysvětlete, jak tuto skutečnost souvisí s tím, že jednotlivé hodinové přírůstky \(\Delta V\) jsou nezáporné.
    \end{tasks}
\end{exercise}


% =========================
% ÚLOHA 20
% =========================
\begin{exercise}[Určení bodů patřících do grafu funkce]
    Mějme graf funkce znázorněný v pravoúhlé soustavě souřadnic.  
    Pomocí něj rozhodněte:

    \begin{center}
    \begin{minipage}[t]{0.45\textwidth}
        \vspace{0pt} % zarovná horní okraje
        \begin{tasks}(1)
            \task Které z~bodů 
            \(A_{1}[1;0]\), \(A_{2}[0;1]\) a \(A_{3}[2;3]\)
            patří do grafu dané funkce?
            \task Vypište všechny body, které patří do grafu této funkce.
        \end{tasks}
    \end{minipage}%
    \hfill
    \begin{minipage}[t]{0.48\textwidth}
        \vspace{0pt}
        \begin{axisgraph}[
            xmin=-3, xmax=4,
            ymin=-2, ymax=3,
            grid=none,
            width=\linewidth,
            height=6cm,
            xlabel={$x$},
            ylabel={$y$},
            minor tick num=1,
            axis lines=middle,
            enlargelimits
        ]
        % --- body grafu ---
        \addplot[
            only marks,
            mark=*,
            mark size=2pt,
            color=black
        ]
        coordinates {
            (-2,1)
            (-1,-1)
            (0,2)
            (1,1)
            (2,0)
            (3,2)
        };
        \end{axisgraph}
    \end{minipage}
    \end{center}
\end{exercise}


% =========================
% ÚLOHA 21
% =========================
\begin{exercise}[Porovnání grafů a předpisů funkcí]
    Přiřaďte ke každé z následujících funkcí odpovídající graf.  
    Všechny grafy jsou zakresleny v pravoúhlé soustavě souřadnic s osami
    \(x\in\langle -1;4\rangle\) a \(y\in\langle -1;6\rangle\).

    \begin{tasks}(2)
        \task \(y = 4x,\quad x \in (0;\infty)\)
        \task \(y = 6x^{2},\quad x \in (0;\infty)\)
        \task \(y = \sqrt{3}\,x,\quad x \in (0;\infty)\)
        \task \(y = x^{3},\quad x \in (0;\infty)\)
    \end{tasks}

    \medskip
    \begin{center}
    % --- čtyři samostatné grafy vedle sebe ---
    \begin{minipage}[t]{0.24\textwidth}
        \vspace{0pt}
        \begin{axisgraph}[
            xmin=-1, xmax=3,
            ymin=-1, ymax=5,
            width=\linewidth,
            height=5cm,
            axis lines=middle,
            grid=none,
            domain=0:3.5,
            samples=100,
        ]
        \addplot[thick, black] {4*x};
        \node at (axis cs:2.5,3) {(1)};
        \end{axisgraph}
    \end{minipage}\hfill
    %
    \begin{minipage}[t]{0.24\textwidth}
        \vspace{0pt}
        \begin{axisgraph}[
            xmin=-1, xmax=3,
            ymin=-1, ymax=5,
            width=\linewidth,
            height=5cm,
            axis lines=middle,
            grid=none,
            domain=0:3.5,
            samples=100,
        ]
        \addplot[thick, black] {6*x^2};
        \node at (axis cs:2.5,3) {(2)};
        \end{axisgraph}
    \end{minipage}\hfill
    %
    \begin{minipage}[t]{0.24\textwidth}
        \vspace{0pt}
        \begin{axisgraph}[
            xmin=-1, xmax=3,
            ymin=-1, ymax=5,
            width=\linewidth,
            height=5cm,
            axis lines=middle,
            grid=none,
            domain=0:3.5,
            samples=100,
        ]
        \addplot[thick, black] {sqrt(3)*x};
        \node at (axis cs:2.5,3) {(3)};
        \end{axisgraph}
    \end{minipage}\hfill
    %
    \begin{minipage}[t]{0.24\textwidth}
        \vspace{0pt}
        \begin{axisgraph}[
            xmin=-1, xmax=3,
            ymin=-1, ymax=5,
            width=\linewidth,
            height=5cm,
            axis lines=middle,
            grid=none,
            domain=0:3.5,
            samples=100,
        ]
        \addplot[thick, black] {x^3};
        \node at (axis cs:2.5,3) {(4)};
        \end{axisgraph}
    \end{minipage}
    % --- konec grafů ---
    \end{center}
\end{exercise}


% =========================
% ÚLOHA 22
% =========================
\begin{exercise}[Tvoření funkce]
    O funkci jsou známy tyto údaje: její definiční obor je interval $\langle-1;4\rangle$,
    obor hodnot je $\langle-2;1\rangle$, funkční hodnoty v bodech $-1;0;4;$ jsou po třadě $0;1;-1;$.
    \begin{tasks}
        \task Kolik existuje funkcí, které splňují všechny tyto podmínky?
        \task Načrtněte graf některé z nich.
    \end{tasks}
\end{exercise}


% =========================
% ÚLOHA 23
% =========================
\begin{exercise}[Tvoření funkce]
    Zkuste nakreslit graf nějaké funkce, pro kterou platí:
    \begin{tasks}
        \task její definiční obor je pětiprvková množina a obor hodnot je množina čtyřprvková;
        \task definiční obor je čtyřprvková množina a obor hodnot je množina pětiprvková. 
    \end{tasks}
\end{exercise}


\setcounter{section}{2}
% =========================
% ÚLOHA 24
% =========================
\begin{exercise}[Tvoření funkce]
    Zkuste nakreslit graf nějaké funkce, pro kterou platí:
    \begin{tasks}
        \task její definiční obor je pětiprvková množina a obor hodnot je množina čtyřprvková;
        \task definiční obor je čtyřprvková množina a obor hodnot je množina pětiprvková. 
    \end{tasks}
\end{exercise}


% =========================
% ÚLOHA 25
% =========================
\begin{exercise}[Vyhodnocování složených funkcí I]
    Jsou dány dvě funkce \(f(x)=3x+1\) a \(g(x)=x^{2}+3x+2\). Vypočítejte:
    \begin{tasks}(3)
        \task \(f(g(1))\)
        \task \(g(f(1))\)
        \task \(f(f(2))\)
        \task \(g(g(-3))\)
        \task \(f(g(a))\)
        \task \(g(f(b))\)
    \end{tasks}
\end{exercise}


% =========================
% ÚLOHA 26
% =========================
\begin{exercise}[Vyhodnocování složených funkcí II]
    Jsou dány dvě funkce \(f(x)=x^{2}+1\) a \(g(x)=(x+1)^{2}\).
    \begin{tasks}(1)
        \task Vypočítejte \(f(g(-2))\).
        \task Vyjádřete předpis funkce \(h(x)=f(g(x))\) a vypočítejte \(h(-2)\).
    \end{tasks}
\end{exercise}


% =========================
% ÚLOHA 27
% =========================
\begin{exercise}[Složené funkce a definiční obory]
    Jsou dány funkce \(f(x)=x-1\), \(g(x)=\sqrt{x}\) a \(h(x)=x+3\). Vypočítejte složené funkce a určete jejich definiční obor:
    \begin{tasks}(2)
        \task \(u_{1}(x)=g(h(f(x)))\)
        \task \(u_{2}(x)=f(h(g(x)))\)
        \task \(u_{3}(x)=f(g(h(x)))\)
        \task \(u_{4}(x)=h(g(f(x)))\)
    \end{tasks}
\end{exercise}


% =========================
% ÚLOHA 28
% =========================
\begin{exercise}[Parametr ve složené funkci]
    Určete hodnotu parametru \(a\in\mathbb{R}\) tak, aby pro funkci \(h(x)=ax+3\) platilo \(h(h(2))=17.\)
\end{exercise}


\setcounter{section}{2}
% =========================
% ÚLOHA 29
% =========================
\begin{exercise}[Rostoucí a klesající části z grafu]
    Jsou dány grafy funkcí \(f\), \(g\), \(h\). Rozhodněte, na kterých intervalech jsou rostoucí a na kterých klesající.
    \begin{center}
    \begin{minipage}[t]{0.32\textwidth}
        \vspace{0pt}
        \begin{axisgraph}[
            xmin=-3, xmax=3,
            ymin=-2, ymax=5,
            width=\linewidth,
            height=5cm,
            axis lines=middle,
            grid=none,
            domain=-3:3,
            samples=200
        ]
        \addplot[thick] {x^2 - 1};
        \node at (rel axis cs:0.8,0.9) {$f$};
        \end{axisgraph}
    \end{minipage}\hfill
    \begin{minipage}[t]{0.32\textwidth}
        \vspace{0pt}
        \begin{axisgraph}[
            xmin=-3, xmax=3,
            ymin=-2, ymax=5,
            width=\linewidth,
            height=5cm,
            axis lines=middle,
            grid=none,
            domain=-3:3,
            samples=200
        ]
        \addplot[thick] {x^5-2*x^3-1/2*x^2+1};
        \node at (rel axis cs:0.9,0.9) {$g$};
        \end{axisgraph}
    \end{minipage}\hfill
    \begin{minipage}[t]{0.32\textwidth}
        \vspace{0pt}
        \begin{axisgraph}[
            xmin=-3, xmax=3,
            ymin=-2, ymax=5,
            width=\linewidth,
            height=5cm,
            axis lines=middle,
            grid=none
        ]
        % lomená čára: rostoucí na (-3,0), klesající na (0,3)
        \addplot[thick] coordinates {(-3,-1) (0,2) (1,1) (3,5)};
        \node at (rel axis cs:0.8,0.9) {$h$};
        \end{axisgraph}
    \end{minipage}
    \end{center}
\end{exercise}


% =========================
% ÚLOHA 30
% =========================
\begin{exercise}[Důkaz rostoucnosti]
    Dokažte, že funkce \(u:\ y=2x-1\) je rostoucí na \(\mathbb{R}\).
\end{exercise}


% =========================
% ÚLOHA 31
% =========================
\begin{exercise}[Důkaz klesající funkce]
    Dokažte, že funkce \(v:\ y=-3x+6\) je klesající na \(\mathbb{R}\).
\end{exercise}


\setcounter{section}{3}
% =========================
% ÚLOHA 32
% =========================
\begin{exercise}[Prosté funkce]
    Rozhodněte, které z následujících funkcí jsou prosté ve svém definičním oboru:
    \begin{tasks}(2)
        \task \(f_{1}:~y=2x+5\)
        \task \(f_{2}:~y=x^{2}-4\)
        \task \(f_{3}:~y=2^{x}\)
        \task \(f_{4}:~y=|x+1|\)
    \end{tasks}
\end{exercise}


\setcounter{section}{4}
% =========================
% ÚLOHA 33
% =========================
\begin{exercise}[Sudé a liché funkce]
    Rozhodněte, které z následujících funkcí \(f_{1}\) až \(f_{12}\) jsou sudé, které liché (v jejich definičním oboru):
    \begin{tasks}(3)
        \task \(f_{1}(x)=x\)
        \task \(f_{2}(x)=\cos x\)
        \task \(f_{3}(x)=\dfrac{1}{x}\)
        \task \(f_{4}(x)=4x\)
        \task \(f_{5}(x)=x^{2}\)
        \task \(f_{6}(x)=\dfrac{4x}{x^{2}-4}\)
        \task \(f_{7}(x)=|x|\)
        \task \(f_{8}(x)=\dfrac{4x}{x^{2}+4}\)
        \task \(f_{9}(x)=x^{3}\)
        \task \(f_{10}(x)=2^{x}\)
        \task \(f_{11}(x)=\dfrac{x^{2}}{|x|+3}\)
        \task \(f_{12}(x)=\log x\)
    \end{tasks}
\end{exercise}


\setcounter{section}{5}
% =========================
% ÚLOHA 34
% =========================
\begin{exercise}[Omezenost funkcí I]
    Rozhodněte, které z následujících funkcí \(g_{1}\) až \(g_{16}\) jsou omezené shora, omezené zdola, resp.\ omezené v definičním oboru:
    \begin{tasks}(4)
        \task \(g_{1}(x)=\cos x\)
        \task \(g_{2}(x)=\sin^{2}x\)
        \task \mbox{$g_{3}(x)=\cos x-1$}
        \task \(g_{4}(x)=2\sin x\)
        \task \(g_{5}(x)=x\)
        \task \(g_{6}(x)=|x|\)
        \task \(g_{7}(x)=-4x\)
        \task \mbox{$g_{8}(x)=-4|x|+3$}
        \task \(g_{9}(x)=x^{2}\)
        \task \(g_{10}(x)=x^{3}\)
        \task \(g_{11}(x)=x^{-1}\)
        \task \(g_{12}(x)=x^{-2}\)
        \task \(g_{13}(x)=2^{x}\)
        \task \(g_{14}(x)=2^{|x|}\)
        \task \(g_{15}(x)=\log x\)
        \task \(g_{16}(x)=-|\log x|\)
    \end{tasks}
\end{exercise}


% =========================
% ÚLOHA 35
% =========================
\begin{exercise}[Omezenost funkcí II – důkazy]
    Dokažte, že následující funkce jsou omezené v jejich definičním oboru:
    \begin{tasks}(1)
        \task \(f_{1}(x)=\dfrac{10}{x^{2}+2}\)
        \task \(f_{2}(x)=x^{2}\), \; pro \(x\in\langle-5;3\rangle\)
        \task \(f_{3}(x)=\dfrac{x}{x+4}\), \; pro \(x\in\mathbb{R}^{+}_{0}\)
    \end{tasks}
\end{exercise}


\setcounter{section}{6}
% =========================
% ÚLOHA 36
% =========================
\begin{exercise}[Extrémy v definičním oboru]
    Rozhodněte, zda funkce \(f_{1}\) až \(f_{9}\) mají v definičním oboru maximum nebo minimum. Pokud ano, určete bod, ve kterém extrém nastává, a jeho hodnotu.
    \begin{tasks}(3)
        \task \(f_{1}(x)=2x\)
        \task \(f_{2}(x)=2x^{2}-1\)
        \task \(f_{3}(x)=\dfrac{x+3}{x}\)
        \task \(f_{4}(x)=|x|+1\)
        \task \(f_{5}(x)=-x^{2}+4\)
        \task \(f_{6}(x)=\dfrac{4x}{x^{2}+4}\)
        \task \(f_{7}(x)=2^{x}+1\)
        \task \(f_{8}(x)=-|\log x|\)
        \task \(f_{9}(x)=\dfrac{1}{x^{2}+2x+1}\)
    \end{tasks}
\end{exercise}


\setcounter{section}{7}
% =========================
% ÚLOHA 37
% =========================
\begin{exercise}[BONUS: Skica grafu podle vlastností]
    Načrtněte graf funkce \(f\), víte-li, že \(D(f)=\langle -3;\infty)-\{0\}\), \(f(-3)=2\); průsečíky s osou \(x\) jsou \(P_{1}[-1;0]\), \(P_{2}[5;0]\); na \(\langle-3;0)\) je \(f\) rostoucí a není omezená shora; na \(\langle-3;3\rangle-\{0\}\) je \(f\) sudá; na \(\langle3;\infty)\) je \(f\) rostoucí a je omezená shora číslem \(h=4\).
    \begin{center}
        \begin{axisgraph}[
            xmin=-6, xmax=8,
            ymin=-4, ymax=6,
            width=12cm,
            height=6.5cm,
            axis lines=middle,
            grid=none,
            xtick={-6,-4,-2,0,2,4,6,8},
            ytick={-4,-2,0,2,4,6},
            minor tick num=1
        ]
        % --- zde načrtněte graf dle vlastností ---
        % \addplot[thick, domain=..., samples=200] { ... };
        \end{axisgraph}
    \end{center}
    \begin{tasks}(1)
        \task Z grafu určete obor hodnot \(H(f)\).
        \task Určete průsečík grafu s osou \(y\) (souřadnice bodu).
        \task Je funkce \(f\) omezená zdola na \(D(f)\)?
        \task Určete maximum funkce \(f\) na \(D(f)\).
        \task Určete minimum funkce \(f\) na \(D(f)\).
        \task Je funkce \(f\) prostá na \(D(f)\)?
        \task Určete alespoň jeden interval, na kterém je \(f\) prostá.
    \end{tasks}

        \medskip
    \textit{Nápověda:}
    \begin{itemize}
        \item Začněte zakreslením bodů \(P_{1}\), \(P_{2}\) a bodu \((-3;2)\).
        \item Využijte, že sudost znamená zrcadlení podle osy \(y\); část grafu vlevo od osy \(y\) tedy určuje i část vpravo.
        \item V bodě \(x=0\) funkce není definována – naznačte přerušení (svislá asymptota).
        \item V intervalu \(\langle3;\infty)\) zakreslete rostoucí křivku, která se shora blíží hodnotě \(y=4\).
    \end{itemize}

\end{exercise}


\end{document}